\documentclass[11pt,a4paper]{article}
\usepackage[T1]{fontenc}
\usepackage[utf8]{inputenc}
\usepackage[italian]{babel}
\usepackage{amsmath,amssymb,amsthm}
\usepackage{booktabs}
\usepackage{geometry}
\geometry{margin=2.5cm}

\theoremstyle{plain}
\newtheorem{proposition}{Proposizione}
\newtheorem{corollary}[proposition]{Corollario}
\theoremstyle{definition}
\newtheorem{remark}[proposition]{Osservazione}

\newcommand{\Qd}{Q_d}
\newcommand{\U}{\mathcal{U}}

\title{Problema 2, Parte 4 (C4)\\[2pt]
\large Cammini in salita su $Q_7$ e $Q_8$: stato del lavoro}
\author{Proof Pursuit}
\date{26 settembre 2026}

\begin{document}
\maketitle

\begin{abstract}
La cella C4 chiede $U(Q_7)$ e $U(Q_8)$. \textbf{Non abbiamo determinato quei
valori}, e questo documento non li dichiara. Riportiamo invece ciò che abbiamo
stabilito: una riduzione esatta per una famiglia naturale di etichettature, la
dimostrazione che quella famiglia \emph{non} contiene l'ottimo già da $d=5$, e i
migliori limiti superiori ottenuti per ricerca locale, con la calibrazione che
ne delimita l'attendibilità. Distinguiamo in modo esplicito ciò che è dimostrato
da ciò che è soltanto verificato al calcolatore su un dominio finito.
\end{abstract}

\section{Definizioni e notazione}

Sia $G$ un grafo semplice finito su $n$ vertici e sia $f\colon V(G)\to\{1,\dots,n\}$
una biiezione (\emph{etichettatura}). Un vertice $v$ è una \emph{valle} se
$f(w)>f(v)$ per ogni vicino $w$ di $v$. Un \emph{cammino in salita} è una
sequenza $(v_1,\dots,v_k)$ con $k\ge1$, in cui $v_1$ è una valle, $v_i$ e
$v_{i+1}$ sono adiacenti e $f(v_1)<\dots<f(v_k)$; una valle isolata conta come
cammino. Sia $U(G)$ il minimo numero di cammini in salita su tutte le
etichettature.

L'ipercubo $\Qd$ ha per vertici $\{0,1\}^d$, con $u\sim v$ quando differiscono
in esattamente una coordinata; $|V(\Qd)|=2^d$ e $|E(\Qd)|=d\,2^{d-1}$. Scriviamo
$|v|$ per il peso di Hamming di $v$.

\paragraph{Convenzione sulle stringhe.} Nelle etichettature riportate, la
coordinata $0$ è il bit meno significativo ed è stampata a destra. La scelta è
irrilevante per il conteggio, per simmetria di $\Qd$, ma va dichiarata.

\paragraph{Ricorrenza di base.} Fissata $f$, orientiamo ogni spigolo dal vertice
di etichetta minore a quello di etichetta maggiore; l'orientazione è aciclica.
Detto $p(v)$ il numero di cammini in salita che terminano in $v$,
\begin{equation}\label{eq:ricorrenza}
p(v)=\begin{cases}1 & \text{se $v$ è una valle},\\[2pt]
\displaystyle\sum_{u\to v}p(u) & \text{altrimenti},\end{cases}
\qquad
\#\{\text{cammini}\}=\sum_{v}p(v).
\end{equation}
Il calcolo è esatto su interi e costa $O(|E|)$, perché l'ordine delle etichette
è già un ordine topologico.

\paragraph{Valori noti.} Il nostro gruppo ha stabilito
\begin{equation}\label{eq:noti}
U(Q_3)=14,\qquad U(Q_4)=34,\qquad U(Q_5)=88,
\end{equation}
con dimostrazione dei limiti inferiori e etichettature esplicite verificate
(celle C1 e C2, consegne \texttt{parte-1.md} e \texttt{parte-2.md}). In questo
documento \eqref{eq:noti} è usata come dato acquisito, non ridimostrata.

\section{Limite inferiore elementare}

\begin{proposition}\label{prop:inferiore}
Per ogni grafo $G$ e ogni etichettatura, il numero di cammini in salita è almeno
$|E(G)|+\nu$, dove $\nu\ge1$ è il numero di valli. In particolare
$U(\Qd)\ge d\,2^{d-1}+1$.
\end{proposition}

\begin{proof}
Ogni vertice $v$ soddisfa $p(v)\ge1$: scendendo da $v$ verso un vicino di
etichetta minore, finché esiste, si raggiunge una valle in un numero finito di
passi, e il cammino percorso all'inverso è in salita e termina in $v$. Se $v$
non è una valle, ciascun suo predecessore $u$ contribuisce $p(u)\ge1$ a
\eqref{eq:ricorrenza}, dunque $p(v)\ge\deg^-(v)$, dove $\deg^-(v)$ conta i
vicini di etichetta minore. Sommando su tutti i vertici e separando le valli,
\[
\sum_v p(v)\;\ge\;\sum_{v\ \text{non valle}}\deg^-(v)\;+\;\nu
\;=\;|E(G)|+\nu,
\]
poiché $\deg^-$ delle valli è nullo e $\sum_v\deg^-(v)=|E(G)|$. L'esistenza di
almeno una valle segue dall'aciclicità dell'orientazione.
\end{proof}

\begin{remark}
Per la griglia $n\times n$ il limite è esatto: $2n^2-2n+1=|E|+1$, che è la
risposta di IMO 2022 P6. Per l'ipercubo \emph{non} lo è: da \eqref{eq:noti},
$U(Q_5)=88$ mentre $|E|+1=81$. L'eccesso è il cuore del problema.
\end{remark}

\section{Etichettature per classi di peso: riduzione esatta}

Chiamiamo \emph{etichettatura per classi di peso} una $f$ ottenuta scegliendo un
ordine delle $d+1$ classi di peso e assegnando le etichette consecutivamente,
classe dopo classe, con ordine interno arbitrario. Sia $\pi(k)$ la posizione
della classe di peso $k$, con $\pi$ permutazione di $\{0,\dots,d\}$.

\begin{proposition}\label{prop:riduzione}
Per un'etichettatura per classi di peso il numero di cammini in salita non
dipende dall'ordine interno alle classi, e vale
\[
N(\pi)\;=\;\sum_{k=0}^{d}\binom{d}{k}P_k,
\]
dove i $P_k$ sono determinati, elaborando le classi in ordine di $\pi$
crescente, da
\[
P_k=
\begin{cases}
1, & \text{se } \pi(k-1)>\pi(k) \text{ e } \pi(k+1)>\pi(k),\\[3pt]
k\,P_{k-1}\,[\pi(k-1)<\pi(k)] \;+\; (d-k)\,P_{k+1}\,[\pi(k+1)<\pi(k)], &
\text{altrimenti},
\end{cases}
\]
con la convenzione che i termini con indice fuori da $\{0,\dots,d\}$ sono omessi.
\end{proposition}

\begin{proof}
\emph{(i) L'orientazione dipende solo da $\pi$.} Se $u\sim v$ in $\Qd$ allora
$u$ e $v$ differiscono in esattamente una coordinata, quindi $\bigl||u|-|v|\bigr|=1$;
in particolare due vertici dello stesso peso non sono mai adiacenti. Per
$u\sim v$ con $|u|\ne|v|$ si ha $f(u)<f(v)$ se e solo se la classe di $u$ precede
quella di $v$, cioè $\pi(|u|)<\pi(|v|)$. L'ordine interno alle classi non
interviene mai, dunque l'orientazione indotta — e con essa il conteggio, che per
\eqref{eq:ricorrenza} dipende solo dall'orientazione — è funzione del solo $\pi$.

\emph{(ii) $p$ è costante sulle classi.} Ogni permutazione $\tau$ delle $d$
coordinate è un automorfismo di $\Qd$ che preserva il peso, quindi per (i)
preserva l'orientazione; dunque $p(\tau v)=p(v)$. Poiché il gruppo delle
permutazioni di coordinate agisce transitivamente su ciascuna classe di peso,
$p(v)$ dipende solo da $|v|$. Poniamo $P_k:=p(v)$ per $|v|=k$.

\emph{(iii) Ricorrenza.} Un vertice $v$ con $|v|=k$ ha esattamente $k$ vicini di
peso $k-1$ (si azzera uno degli uni) e $d-k$ vicini di peso $k+1$ (si alza uno
degli zeri). Per (i) i primi sono predecessori se e solo se $\pi(k-1)<\pi(k)$, i
secondi se e solo se $\pi(k+1)<\pi(k)$. Il vertice è una valle esattamente
quando nessuna delle due condizioni vale. Sostituendo in \eqref{eq:ricorrenza} e
usando (ii) si ottiene la formula. Sommando su tutti i vertici e raggruppando
per peso si ha $N(\pi)=\sum_k\binom{d}{k}P_k$. L'elaborazione in ordine di $\pi$
crescente è lecita perché i predecessori di una classe appartengono a classi che
la precedono.
\end{proof}

\begin{corollary}\label{cor:enumerazione}
Il minimo di $N(\pi)$ sulle etichettature per classi di peso si calcola
esattamente enumerando le $(d+1)!$ permutazioni, con costo $O(d)$ ciascuna e
aritmetica intera.
\end{corollary}

La Tabella~\ref{tab:classi} riporta l'esito dell'enumerazione completa.

\begin{table}[h]
\centering
\begin{tabular}{crrrl}
\toprule
$d$ & $|E|+1$ & $\min_\pi N(\pi)$ & eccesso su $|E|$ & confronto con $U(\Qd)$\\
\midrule
3 & 13   & 14   & 2   & $=U(Q_3)$\\
4 & 33   & 34   & 2   & $=U(Q_4)$\\
5 & 81   & 92   & 12  & $>U(Q_5)=88$\\
6 & 193  & 214  & 22  & ignoto\\
7 & 449  & 506  & 58  & ignoto\\
8 & 1025 & 1138 & 114 & ignoto\\
\bottomrule
\end{tabular}
\caption{Minimo esatto entro le etichettature per classi di peso.
Enumerazione completa delle $(d+1)!$ permutazioni, $34{,}5$~s complessivi per
$d\le8$. Ogni riga è stata ricontrollata con un conteggio diretto sul grafo,
che non usa la Proposizione~\ref{prop:riduzione}.}
\label{tab:classi}
\end{table}

\begin{proposition}\label{prop:nonottimale}
Le etichettature per classi di peso non sono ottimali per $d=5$: il loro minimo
vale $92$, mentre $U(Q_5)=88$.
\end{proposition}

\begin{proof}
Il valore $92$ è il minimo di $N(\pi)$ sulle $6!=720$ permutazioni, ottenuto per
enumerazione completa con aritmetica esatta (Corollario~\ref{cor:enumerazione});
$U(Q_5)=88$ è il risultato della cella C2. Poiché $92>88$, nessuna etichettatura
per classi di peso realizza il minimo.
\end{proof}

\begin{remark}[Conseguenza operativa]
La Proposizione~\ref{prop:nonottimale} è il motivo per cui \textbf{non}
consegniamo $506$ e $1138$ come valori di $U(Q_7)$ e $U(Q_8)$. La famiglia
riproduce i valori corretti per $d\le4$ e si rivela subottimale a $d=5$: non c'è
ragione di aspettarsi che ridiventi ottimale per $d\ge6$. Essendo C4 valutata sul
valore, una risposta plausibile ma non dimostrata varrebbe zero.
\end{remark}

\section{Una costruzione esatta da insiemi che spezzano i cicli}

La famiglia per classi di peso fallisce perché produce troppi pochi vertici con
$p=1$. La costruzione seguente massimizza invece proprio quelli.

\begin{proposition}\label{prop:costruzione}
Sia $I\subseteq V(\Qd)$ un insieme \emph{indipendente} tale che $\Qd-I$ sia una
\emph{foresta}. Allora esiste un'etichettatura di $\Qd$ con esattamente
\[
2^d+(d-1)\,|I|
\]
cammini in salita. In particolare
$U(\Qd)\le 2^d+(d-1)\,\nabla_{\mathrm{ind}}(\Qd)$, dove
$\nabla_{\mathrm{ind}}(\Qd)$ è la minima cardinalità di un tale $I$.
\end{proposition}

\begin{proof}
Costruiamo $f$ così: assegniamo le etichette $1,\dots,2^d-|I|$ ai vertici di
$\Qd-I$, visitando ogni albero della foresta in ampiezza a partire da una sua
radice scelta arbitrariamente, in modo che ogni vertice riceva etichetta maggiore
del proprio padre; assegniamo poi le etichette restanti, tutte maggiori, ai
vertici di $I$ in ordine arbitrario.

Sia $v\notin I$. I suoi vicini sono di due tipi. Quelli in $I$ hanno etichetta
maggiore per costruzione, quindi non sono predecessori. Quelli in $\Qd-I$ sono i
vicini di $v$ nella foresta: sono il padre di $v$ (etichetta minore) e i figli di
$v$ (etichetta maggiore), e il padre esiste se e solo se $v$ non è radice.
Dunque $\deg^-(v)=0$ se $v$ è radice e $\deg^-(v)=1$ altrimenti. Per
\eqref{eq:ricorrenza}, per induzione lungo l'albero a partire dalla radice,
$p(v)=1$ per ogni $v\notin I$: le radici sono valli, e ogni altro vertice eredita
il valore $1$ dal solo padre.

Sia ora $v\in I$. Poiché $I$ è indipendente, tutti i $d$ vicini di $v$ stanno in
$\Qd-I$ e hanno quindi etichetta minore. Per \eqref{eq:ricorrenza},
$p(v)=\sum_{u\sim v}p(u)=d\cdot 1=d$.

Sommando, $\sum_v p(v)=(2^d-|I|)\cdot 1+|I|\cdot d=2^d+(d-1)|I|$.
\end{proof}

\begin{remark}
La Proposizione~\ref{prop:costruzione} riconduce un problema di minimo su $2^d!$
etichettature a un problema di minimo su sottoinsiemi di vertici: trovare il più
piccolo insieme indipendente il cui complemento sia aciclico. È una variante
\emph{indipendente} del problema del \emph{decycling set} (insieme di vertici di
retroazione).
\end{remark}

\section{Risultati sperimentali}

Abbiamo attaccato il problema con due procedure \textbf{indipendenti nel metodo},
non due riscritture della stessa idea:

\begin{enumerate}
\item \textbf{Ricottura sulle etichettature} (\texttt{c4\_ricerca\_locale.py}):
stato $=$ permutazione dei $2^d$ vertici, mosse $=$ scambio e spostamento di
etichette, punteggio $=$ conteggio esatto via \eqref{eq:ricorrenza}. Questa
procedura \textbf{non conosce} la Proposizione~\ref{prop:costruzione}.
\item \textbf{Ricottura sugli insiemi} (\texttt{c4\_insieme\_decycling.py}):
stato $=$ sottoinsieme $I$ di vertici, costo $=|I|$ più penalità per coppie
adiacenti in $I$ e per spigoli in eccesso rispetto a una foresta; il valore si
ottiene poi dalla Proposizione~\ref{prop:costruzione}.
\end{enumerate}

\paragraph{Calibrazione.} Entrambe le procedure sono state prima eseguite su
$d=5$, dove la risposta $88$ è nota e dimostrata. Entrambe l'hanno ritrovata: la
prima in $45$~s ($281\,600$ mosse), la seconda in $50$~s (trovando $|I|=14$).
Una procedura che non ritrova un valore noto non è credibile su valori ignoti.

\paragraph{Stabilità.} Su $d=6$ la ricottura sulle etichettature è stata
ripetuta con tre semi indipendenti ($1$, $7$, $13$), per $500$--$600$~s
ciascuna: tutte e tre si fermano su $204$, mai sotto. La ricerca sugli insiemi,
partendo da tutt'altro spazio di stati, trova $|I|=28$ e quindi di nuovo $204$.

\begin{table}[h]
\centering
\begin{tabular}{crrrrl}
\toprule
$d$ & $|I|$ & $2^d+(d-1)|I|$ & ricottura etichettature & stato\\
\midrule
3 & 3   & 14   & ---   & $=U(Q_3)$, dimostrato (C1)\\
4 & 6   & 34   & ---   & $=U(Q_4)$, dimostrato (C1)\\
5 & 14  & 88   & $88$  & $=U(Q_5)$, dimostrato (C2)\\
6 & 28  & 204  & $204$ & limite superiore\\
7 & 56  & 464  & $464$ & limite superiore\\
8 & 112 & 1040 & $1040$& limite superiore\\
\bottomrule
\end{tabular}
\caption{Le due procedure concordano su ogni riga in cui entrambe sono state
eseguite. Ogni etichettatura prodotta è stata ricontata con un conteggio
indipendente, e ogni $I$ è stato verificato indipendente e con complemento
aciclico.}
\label{tab:risultati}
\end{table}

\begin{proposition}\label{prop:bound}
$U(Q_6)\le 204$, \quad $U(Q_7)\le 464$, \quad $U(Q_8)\le 1040$.
\end{proposition}

\begin{proof}
Per ciascun $d\in\{6,7,8\}$ esibiamo nei file allegati un insieme $I$ di
cardinalità $28$, $56$, $112$ rispettivamente, verificato indipendente e con
$\Qd-I$ foresta, e l'etichettatura esplicita che ne discende. La
Proposizione~\ref{prop:costruzione} dà il valore; il conteggio diretto sulle
etichettature lo conferma.
\end{proof}

\begin{remark}[Congettura]
I valori della Tabella~\ref{tab:risultati} soddisfano tutti
$U(\Qd)=2^d+(d-1)\nabla_{\mathrm{ind}}(\Qd)$, e per $d=3,4,5$ questo coincide con
il valore \emph{dimostrato}. Congetturiamo quindi
\[
U(Q_7)=464,\qquad U(Q_8)=1040,
\]
ma \textbf{non lo dimostriamo}: manca il limite inferiore, cioè la prova che
nessun insieme indipendente più piccolo renda $\Qd$ una foresta e che nessuna
etichettatura fuori da questa costruzione faccia meglio.
\end{remark}

\begin{remark}[Un vincolo su $d=9$]
La successione $|I|=3,6,14,28,56,112$ raddoppia da $d=5$ in poi. Se il
raddoppio proseguisse si avrebbe $\nabla_{\mathrm{ind}}(Q_9)\le224$, e la
Proposizione~\ref{prop:costruzione} darebbe $U(Q_9)\le 512+8\cdot224=2304$. Ma
gli organizzatori dichiarano $U(Q_9)\ge2368$. Poiché la
Proposizione~\ref{prop:costruzione} è dimostrata, le due cose sono
incompatibili: assumendo corretto il limite degli organizzatori,
\[
\nabla_{\mathrm{ind}}(Q_9)\;\ge\;232,
\]
cioè il raddoppio \emph{deve} rompersi a $d=9$. È un'informazione utile per C5 e
C6: indica che la costruzione va modificata, non estesa per inerzia.
\end{remark}

\section{Che cosa consideriamo risolto}

\begin{itemize}
\item \textbf{Dimostrato:} la Proposizione~\ref{prop:inferiore}
(limite inferiore $|E|+\nu$) e la Proposizione~\ref{prop:riduzione}
(riduzione esatta per classi di peso).
\item \textbf{Verificato al calcolatore su dominio finito:} la
Tabella~\ref{tab:classi}, per enumerazione completa delle $(d+1)!$ permutazioni
con $d\le8$, in aritmetica intera, con controllo incrociato indipendente.
\item \textbf{Dato acquisito da altra cella:} i valori \eqref{eq:noti}.
\item \textbf{Non stabilito:} $U(Q_7)$ e $U(Q_8)$. La cella C4 resta
\textbf{non risolta}; quanto sopra è progresso parziale.
\end{itemize}

\end{document}
